% file: sections/intro.tex

\section{研究背景}

%%%%%%%%%%%%%%%
\begin{frame}{}
  \begin{center}
    {\large Abstract Data Types} \ncite{Liskov:VHLL74} \\[8pt]

    (单线程; 顺序语义)
  \end{center}

  \begin{columns}
    \column{0.40\textwidth}
      \fignocaption{width = 0.60\textwidth}{figs/alg-ds-wirth}
    \pause
    \column{0.40\textwidth}
      \fignocaption{width = 0.25\textwidth}{figs/adt}
  \end{columns}
\end{frame}
%%%%%%%%%%%%%%%

%%%%%%%%%%%%%%%
\begin{frame}{}
  \begin{center}
    {\large Concurrent Data Types} \ncite{Herlihy:TOPLAS90} \\[8pt]

    (多线程; 并发语义)
  \end{center}

  \begin{columns}
    \column{0.40\textwidth}
      \fignocaption{width = 0.70\textwidth}{figs/taomp-herlihy}
    \pause
    \column{0.50\textwidth}
      \fignocaption{width = 0.45\textwidth}{figs/cdt}

      \only<3->{
	\begin{center}
	  \hl{\Large PL} {(Programming Language)}
	\end{center}
      }
  \end{columns}
\end{frame}
%%%%%%%%%%%%%%%

%%%%%%%%%%%%%%%
\begin{frame}{}
  \begin{center}
    \hl{\large Replicated Data Types} \ncite{Shapiro:TR11} \ncite{Burckhardt:POPL14} \\[8pt]

    (多副本; 复制语义)
  \end{center}

  \pause
  \fignocaption{width = 0.30\textwidth}{figs/rdt}

  \pause
  \begin{center}
    \hl{\Large DC} {(Distributed Computing)}
  \end{center}
\end{frame}
%%%%%%%%%%%%%%%

%%%%%%%%%%%%%%%
% \begin{frame}{}
%   \begin{columns}
%     \column{0.50\textwidth}
%       \fignocaption{width = 0.60\textwidth}{figs/keep-calm-why-bother}
%     \pause
%     \column{0.50\textwidth}
%       \centerline{\hl{\red{\Huge 新平台}}}
%   \end{columns}
% \end{frame}
%%%%%%%%%%%%%%%

%%%%%%%%%%%%%%%
\begin{frame}{}
  \begin{center}
    {\hl{\Large 新平台: 大规模分布式系统}}
  \end{center}

  \fignocaption{width = 0.50\textwidth}{figs/sina-weibo-world-map.pdf}

  \pause
  \vspace{0.10cm}
  \begin{center}
    低延迟 \quad 高可用性 {\small (4个9)} \quad 高容错性 \quad 高可扩展性
  \end{center}

  % \begin{columns}
  %   \column{0.40\textwidth}
  %   新浪微博社交应用~\footnotemark:
  %   \begin{itemize}
  %     \item 日均用户近一亿名
  %     \item 日均消息近一亿条
  %   \end{itemize}
  %   \pause
  %   \column{0.50\textwidth}
  %   特性需求: 
  %   \begin{itemize}
  %     \item 低延迟, 高可用性 (4个9~\footnotemark)
  %     \item 高容错性, 高可扩展性
  %   \end{itemize}
  % \end{columns}
  % 
  % \footnotetext[1]{2015第三季度; 数据来自 \href{https://www.chinainternetwatch.com/15740/weibo-q3-2015/}{China Internet Watch}.}
  % % See http://tex.stackexchange.com/a/340079/23098
  % \alt<1>{\let\thefootnote\relax\footnotetext{~}}{\footnotetext[2]{数据来自 \href{http://www.infoq.com/cn/articles/weibo-platform-availability-9999}{InfoQ}.}}
\end{frame}
%%%%%%%%%%%%%%%
\begin{frame}{}
  \graphicspath{{tikz/}}
  \begin{figure}[h!]
    \centering
    \begin{adjustbox}{max totalsize = {0.50\textwidth}{1.00\textheight}, center}
      % File: beamer overlay for distributed-data.tex
% Created: July 23, 2016

\begin{tikzpicture}[rc/.style = {circle, minimum size = 4cm, fill = red},
  br/.style = {rectangle, rounded corners, minimum size = 4cm, fill = blue},
  gs/.style = {shape = star, scale = 8, star points = 5, star point ratio = 1.65, fill = green},
  dash pattern = {on 20pt off 10pt},
  partition/.style = {draw, lightgray, line width = 8pt},
  replication/.style = {draw, dash phase = 0pt, line width = 10pt, #1},
  access/.style = {draw, > = stealth, <->, dash phase = 0pt, line width = 6pt, brown}]

  \uncover<1->{
  \begin{pgfonlayer}{background}
	\node[opacity = 0.80] (map) {\includegraphics[scale = 0.50]{figs/world-map-logo.jpg}};
  \end{pgfonlayer}

  % shared data
  \begin{scope}[yshift = -18cm]
	\node[rc] (circ) {};
	\node[br, left = 2cm of circ] (rect) {};
	\node[gs, right = 2cm of circ] (st) {};

	\begin{pgfonlayer}{background}
	  \node[rectangle, rounded corners, fill = black!30, fit = (circ) (rect) (st), inner sep = 5pt, 
		label = {[font = \fontsize{65}{70}\sffamily\bfseries] below : Application Data}] (data) {};
	\end{pgfonlayer}
  \end{scope}
  }

  \uncover<2->{
  % partition
  \node[rc] (circ1) at (20, 10) {};
  \node[br] (rect1) at (-25, 10) {};
  \node[gs] (st1) at (28, -12) {};

  \draw[partition] (circ) to (circ1);
  \draw[partition] (rect) to (rect1);
  \draw[partition] (st) to (st1);
  }

  \uncover<3->{
  % replication
  \node[rc] (circ2) at (-15, -10) {};
  \node[rc] (circ3) at (-10, 18) {};

  \node[br] (rect2) at (10, 12) {};

  \node[gs] (st2) at (0,0) {};
  \node[gs] (st3) at (20, 5) {};

  \path[replication = {red}] (circ1) to (circ2) 
  	(circ2) to (circ3)
	(circ3) to (circ1);

  \path[replication = {blue}] (rect1) to (rect2);

  \path[replication = {green}] (st1) to (st2) 
  	(st2) to (st3)
	(st3) to (st1);
  }

%  \uncover<5->{
%  % sina-weibo clients
%  \def\SCALE{0.15}
%  \node[] (sina1) at (-25, 0) {\includegraphics[scale = \SCALE]{figs/sina-weibo.png}};
%  \node[] (sina2) at (10, 5) {\includegraphics[scale = \SCALE]{figs/sina-weibo.png}};
%  \node[] (sina3) at (5, -10) {\includegraphics[scale = \SCALE]{figs/sina-weibo.png}};
%
%  % access
%  \path[every edge/.append style = {access}] (sina1) edge (rect1)
%  				(sina1) edge (circ2)
%				(sina2) edge (rect2)
%				(sina2) edge (st2)
%				(sina3) edge (circ1)
%				(sina3) edge (st2);
%  }
\end{tikzpicture}

    \end{adjustbox}
  \end{figure}

  \vspace{0.20cm}
  \begin{center}
    \begin{minipage}{0.65\textwidth}
      \red{\large 分布数据 \term{distributed data}:}

      \vspace{0.20cm}
      \begin{enumerate}
	\item<2-> 分区 \term{partition}: 水平扩展
	\item<3-> \hl{副本 \term{replication}}: 就近访问, 容灾备份
      \end{enumerate}
    \end{minipage}
  \end{center}
\end{frame}
%%%%%%%%%%%%%%%

%%%%%%%%%%%%%%%
\begin{frame}{}
  \begin{center}
    \begin{minipage}{0.50\textwidth}
      {\large 复制数据类型} \ncite{Shapiro:TR11} 

      \vspace{0.20cm}
      \begin{itemize}
	\setlength{\itemsep}{4pt}
	\item Read/Write Register
	\item Counter
	\item Set
	\item List
	\item HashMap
	\item Disjoint Set
	\item Graph
	\item $\dots$
      \end{itemize}
    \end{minipage}
  \end{center}
\end{frame}
%%%%%%%%%%%%%%%

%%%%%%%%%%%%%%%
\begin{frame}{}
  \begin{columns}
    \column{0.40\textwidth}
      \fignocaption{width = 0.98\textwidth}{figs/what-is-new}
    \column{0.60\textwidth}
      \centerline{\hl{\red{\Huge 新问题, 新挑战}}}
  \end{columns}

  \pause
  \vspace{0.30cm}
  \fignocaption{frame, width = 0.80\textwidth}{figs/rdt-popl}
  \centerline{\ncite{Burckhardt:POPL14}}
\end{frame}
%%%%%%%%%%%%%%%

%%%%%%%%%%%%%%%
\begin{frame}{}
  \begin{center}
    \resizebox{0.65\textwidth}{!}{% file: rdt-research-two-work.tex

\begin{tikzpicture}[
    root concept/.append style = {concept color = teal!40, sibling angle = 90, font = \LARGE, minimum size = 0pt, red},
    level 1 concept/.append style = {concept color = teal!40, sibling angle = 90, font = \large, violet},
    level 2 concept/.append style = {concept color = teal!40, font = \large, purple},
    every annotation/.style = {fill = red!20, font = \sf, align = center, font = \large,
      minimum size = 5pt, inner sep = 6pt, text width = 2.6cm}
  ] 
  \path[mindmap]
    node (spec) [concept] {\textsl{Specification}}
    [counterclockwise from = 225]

    child[] {
      node (impl) [concept] {\textsl{Imple\-mentation}}
      [counterclockwise from = -90]
      child { node (jupiter) [concept, font = \large] {Jupiter} }
    }
    child[] {
      node (veri) [concept] {\textsl{Veri\-fication}}
      [clockwise from = -90]
      child { node (vpc) [concept, font = \large] {VPC} }
    };

    \uncover<2->{\node [annotation, below = 0.20cm of jupiter.south, text width = 1.0cm] {List};}
    \uncover<3->{\node [annotation, below = 0.10cm of vpc.south] {R/W Register};}
    \uncover<4->{\node [annotation, fill = yellow!20, above = 0.40cm of veri.center, text width = 1.6cm] {Checking};}

    \node [annotation, fill = brown!20, below = 1.00cm of spec.south, font = \Huge, inner sep = 12pt, brown] {RDT};
\end{tikzpicture}
}
  \end{center}
\end{frame}
%%%%%%%%%%%%%%%
